\documentclass[11pt]{article}
\usepackage{booktabs,xcolor,colortbl,caption,array,graphicx,amsmath}
\usepackage[margin=1.5cm]{geometry}
\begin{document}\small\setcounter{table}{6}
\begin{table}[!h]
\centering
\caption{\label{tab:tab:summ_scen}NFA et dark matter sous deux scénarios, moy. 2013--2022 (\% du PIB)}
\centering
\fontsize{10.5}{12.5}\selectfont
\begin{tabular}[t]{llrrrrrr}
\toprule
\multicolumn{2}{c}{\textbf{ }} & \multicolumn{3}{c}{\textbf{Mesures NFA}} & \multicolumn{3}{c}{\textbf{Dark matter}} \\
\cmidrule(l{3pt}r{3pt}){3-5} \cmidrule(l{3pt}r{3pt}){6-8}
Entité & Groupe & NIIP & NFA$_A$ & NFA$_B$ & DM$_A$ & DM$_B$ & B$-$A\\
\midrule
Chine & emerging & 14.4 & -10.0 & -8.7 & -24.4 & -23.1 & 1.2\\
Union européenne & advanced & -8.5 & -4.5 & -4.5 & 3.9 & 3.9 & 0.0\\
Japon & advanced & 63.0 & 73.5 & 73.5 & 10.5 & 10.5 & 0.0\\
États-Unis & privilege & -52.7 & 11.9 & 6.1 & 64.7 & 58.8 & -5.9\\
\bottomrule
\multicolumn{8}{l}{\rule{0pt}{1em}\textit{\textit{Notes:} }}\\
\multicolumn{8}{l}{\rule{0pt}{1em}Scén. A: taux G\&R (2006) universels. Scén. B: GRG (2017). Privilège (USA, CHE, GBR): r$_{\text{FDI}}$=9\%, r$_{\text{dette}}$=2\%. Émergents (CHN): r$_{\text{FDI}}$=6\%, r$_{\text{dette}}$=5\%. Un taux plus élevé réduit le stock capitalisé pour le même flux de revenu. UE: pays avancés → Scén. B = Scén. A.}\\
\end{tabular}
\end{table}
\end{document}
